% Le triangle du feu (L'essentiel, chapitre 10)
\begin{tikzpicture}[x=1cm,y=1cm,
  sommet/.style={draw=otNavy,fill=otNavyTint,line width=0.6pt,rounded corners=1.4mm,align=center,
    font=\otMedium\fontsize{8.5}{10}\selectfont,text=otNavy,inner sep=4pt}]
\draw[otNavy!50,line width=1.2pt,line join=round] (0,0) -- (5,0) -- (2.5,3.6) -- cycle;
\fill[otOrange,opacity=0.9] (2.5,0.55) .. controls (1.85,0.7) and (1.9,1.5) .. (2.3,1.95) .. controls (2.45,2.15) and (2.5,2.35) .. (2.55,2.6)
  .. controls (2.8,2.2) and (3.15,1.7) .. (3.05,1.2) .. controls (3.0,0.8) and (2.8,0.58) .. (2.5,0.55);
\fill[yellow!70!white] (2.5,0.65) .. controls (2.2,0.8) and (2.2,1.4) .. (2.5,1.8) .. controls (2.75,1.4) and (2.8,0.85) .. (2.5,0.65);
\node[sommet] at (0,-0.15) {combustible\\(réducteur)};
\node[sommet] at (5,-0.15) {comburant : dioxygène\\(oxydant)};
\node[sommet] at (2.5,3.85) {énergie d'activation\\(flamme, étincelle)};
\end{tikzpicture}
\par\smallskip\legende{Une combustion nécessite les trois éléments du triangle du feu ; en supprimer un l'éteint.}
